\section*{Limitations}
This study evaluates a bounded set of hypotheses, not the proposition that
all graph-conditioned credit maintenance is ineffective or effective.
Earlier candidate decisions are retained, including unsuccessful routes.
The current design uses one task, two model checkpoints, a small fixed
graph menu, and a single exogenous topology change per episode. It does not
evaluate long trajectories with repeated graph edits, adversarial drift,
model replacement, or generalization to a second task. A terminal decision
that reverses an initial deletion preference is a restoration proxy, not
an observed delete-and-restore deployment trajectory.
The service graph is selected once per episode after a bounded calibration
window. The periodic baseline schedules feedback within that window; it
is not a continuously interleaved production refresh policy.

All active nodes in the current candidate receive the same full raw
document context. This isolates communication and execution-cost effects
within that access regime; it does not represent agents holding mutually
private evidence. Earlier partial-context experiments are separate cohorts
and cannot establish a causal advantage for the common-context regime.
The model weights may have encountered benchmark material during
pretraining, which our exclusion ledger cannot establish or rule out.

The Dirichlet model is an intentionally simple decision statistic. It
retains old counts and does not infer the true amount of graph-induced
distribution shift. Its marginal edge distributions do not model every
dependence arising from shared questions and baseline executions. The
one-step expected value calculation is not a proof of optimal multi-step
sampling. A graph-topology signal can exist while providing no additional
decision value beyond a static protection rule.

Gold benchmark scoring supplies feedback in replay. Human labeling time,
human financial cost, and production feedback availability were not
measured. Token utility uses a declared research tradeoff, not a monetary
conversion or a common physical price for different model sizes. Kernel
time and GPU residence are separately reported; neither establishes
latency, energy, or monetary savings in an online system. Exact-request
caching reduces data-collection work, while the policy ledger separately
charges the requests each policy would purchase. This is counterfactual
replay, not a live policy deployment.

The common token cap is an upper bound and may be nonbinding. Equal caps
must not be described as equal realized spending. No-initialization static
controls deliberately test whether paying for maintenance is necessary.
Service-prefix curves retain the originally fixed horizon and do not
estimate policies optimized for every shorter horizon.

Repeated generation seeds share question groups and are not additional
independent data samples. We report question support, paired point
differences, correction and breakage counts, and episode and seed
breakdowns. The new experiment uses no confidence intervals or bootstrap;
its frozen numerical gates are resource-allocation criteria, not tests of
statistical significance or population noninferiority. Previous exposed
data and post-hoc analyses cannot supply independent confirmation.

The literature review records whether claims were checked in an abstract
or in specific paper sections. Published graph-optimization systems were
not all independently reimplemented in this study. Comparisons with
internal baselines therefore cannot establish superiority over the full
external state of the art. Computational replays and a separate arithmetic
implementation establish limited verification scope, not independent
human peer review or submission readiness.
